\SourceLeaf{360}%
and therefore forgoes the advantages that other methods can draw from the shorter distance, a far more common case in practice. The useful application of \textsc{Bessel}'s method will therefore be limited to cases in which the two points are connected through a longer chain of triangles instead of directly by a side of an actually measured triangle. But then one must rightly ask. how the \textit{data} for the problem are to be obtained, namely the actual length of the geodesic line joining the two points and the angle it makes with the meridian at one endpoint? To determine these by a purely \textit{spherical} calculation of the intervening triangles, as \textsc{Bessel} did for the Prussian degree measurement of limited extent, would no longer be admissible at a much greater distance. But if this passage is calculated on the spheroid, this alone will require as much work as determining at once the latitude, longitude and reverse azimuth of each succeeding point. These remarks also apply to \textsc{Ivory}'s method of solution, which closely resembles \textsc{Bessel}'s but gives little attention to what is actually needed in practice.

A few further remarks may be made here about the method given in the present memoir.

The formulas give directly the \textit{differences} between the two latitudes and between the two azimuths, as well as the difference in longitude. The smallness of these differences is what makes the calculation accurate without requiring logarithms with more digits, provided that for the azimuths one is reckoned from the south side of the meridian and the other from the north side. The symmetry and simplicity of the formulas, on the other hand, come from their depending in the first instance on the mean of the two latitudes and the mean of the two azimuths, instead of on the latitude and azimuth at one endpoint. It follows that the formulas can be used to solve the problem as stated above only by an indirect procedure, or by successive approximation. Skilled calculators who are familiar with the devices used in such operations will hardly find this inconvenient, especially since it can be assumed that close approximations to the quantities to be determined are almost always already available by the time the precise %
